% Source: https://artofproblemsolving.com/community/c3h2151089_rational_root_theorem
% Stable identity: aops-topic-2151089-q01
Let $f(x)=x^3+3x^2+4x-11$.

\textbf{(a)} By the rational root theorem, a rational root $\dfrac {p}{q}$ in lowest terms must have $p\mid11$ and $q\mid1$, so it is one of $\pm1,\pm11$. But
\[f(1)=-3,\quad f(-1)=-13,\quad f(11)=1727,\quad f(-11)=-1023,\]
so there is no rational root, and in particular $\dfrac{121}{100}$ is not a root.

\textbf{(b)} Substitute $x=y-1$:
\[f(y-1)=(y-1)^3+3(y-1)^2+4(y-1)-11=y^3+y-13 .\]
The function $y^3+y$ is strictly increasing, so $f$ is strictly increasing and has at most one real root; it has one, since it changes sign. Moreover
\[f(1.2)=1.728+4.32+4.8-11=-0.152<0,\qquad f(1.21)=1.771561+4.3923+4.84-11=0.003861>0,\]
so the unique real root lies strictly between $1.2$ and $1.21$ (it is about $1.2097$). The calculator displays this irrational number rounded to two decimal places; $1.21$ is only an approximation, not an exact root.
